% Partie : intuitionistic-logic
% Chapitre : introduction
% Section : axiomatic-derivations

\documentclass[../../../include/open-logic-section]{subfiles}

\begin{document}

\olfileid[fr]{int}{int}{axd}

\olsection{\usetoken{P}{derivation} axiomatiques}

Les !!{derivation}s axiomatiques pour la logique propositionnelle
intuitionniste constituent les systèmes de !!{derivation}
conceptuellement les plus simples et historiquement les premiers.
Leur fonctionnement est le même qu'en logique propositionnelle classique.

\begin{defn}[!!^{derivability}]
Soit $\Gamma$ un ensemble de !!{formula}s du langage $\Lang L$.
Une \emph{!!{derivation}} à partir de $\Gamma$ est une suite finie
$!A_1$, \dots,~$!A_n$ de !!{formula}s telle que, pour chaque
$i \le n$, l'une des conditions suivantes soit remplie :
\begin{enumerate}
\item $!A_i \in \Gamma$ ; ou
\item $!A_i$ est un axiome ; ou
\item $!A_i$ découle par modus ponens de formules précédentes
  $!A_j$ et $!A_k$, avec $j < i$ et $k < i$ ; autrement dit,
  $!A_k \ident !A_j \lif !A_i$.
\end{enumerate}
\end{defn}

\begin{defn}[Axiomes]
L'ensemble $\PAx$ des \emph{axiomes} de la logique propositionnelle
intuitionniste comprend toutes les !!{formula}s des formes suivantes :
\begin{align}
  & (!A \land !B) \lif !A \ollabel{ax:land1}\\
  & (!A \land !B) \lif !B \ollabel{ax:land2}\\
  & !A \lif (!B \lif (!A \land !B)) \ollabel{ax:land3}\\
  & !A \lif (!A \lor !B) \ollabel{ax:lor1}\\
  & !A \lif (!B \lor !A) \ollabel{ax:lor2}\\
  & (!A \lif !C) \lif ((!B \lif !C) \lif ((!A \lor !B) \lif !C)) \ollabel{ax:lor3}\\
  & !A \lif (!B \lif !A) \ollabel{ax:lif1}\\
  & (!A \lif (!B \lif !C)) \lif ((!A \lif !B) \lif (!A \lif !C)) \ollabel{ax:lif2}\\
  & \lfalse \lif !A \ollabel{ax:lfalse1}
\end{align}
\end{defn}

\begin{defn}[!!^{derivability}]
Une !!{formula}~$!A$ est \emph{!!{derivable}} à partir de
$\Gamma$, ce que l'on note $\Gamma \Proves !A$, s'il existe une
!!{derivation} à partir de $\Gamma$ dont la dernière formule est $!A$.
\end{defn}

\begin{defn}[Théorèmes]
Une !!{formula}~$!A$ est un \emph{théorème} s'il existe une
!!{derivation} de~$!A$ à partir de l'ensemble vide. Nous écrivons
$\Proves !A$ si $!A$ est un théorème, et $\Proves/ !A$ sinon.
\end{defn}

\begin{prop}
  Si $\Gamma \Proves !A$ dans ce système de !!{derivation}
  axiomatique intuitionniste, alors $\Gamma \Proves !A$ aussi
  en logique classique. En particulier, tout théorème
  intuitionniste $!A$ est également un théorème classique.
\end{prop}

\begin{proof}
  Chacun des axiomes intuitionnistes figure aussi parmi les
  axiomes du système classique considéré ici ; la règle de
  modus ponens est la même. Toute !!{derivation} intuitionniste
  est donc aussi une !!{derivation} classique.
\end{proof}

\end{document}
